\paragraph{Pick--Poisson 负谱体积势的二次 coercive 几何与 Schur 瓶颈}
正文已经把 Pick--Poisson 负谱体积势写成端点自能与伪双曲互能之和，并给出并集交叉互能恒等式、固定总端点通量下的极值上界、以及 Schur 补账本的逐步外部化记账。把这些接口合并后，可把 \(-\log\det \mathsf P\) 进一步压缩为一条真正的二次 coercive 几何链。

\leanverified{paper\_derived\_pick\_poisson\_clustered\_quadratic\_coercive\_lower\_bound}
\begin{theorem}[分簇缺陷能量的二次强制下界]\label{thm:derived-pick-poisson-clustered-quadratic-coercive-lower-bound}
设有限缺陷多重集分解为互不相交并集
\[
A=\bigsqcup_{r=1}^{R}A_r,
\qquad
\kappa_r:=|A_r|,
\qquad
\kappa:=\sum_{r=1}^{R}\kappa_r.
\]
若存在常数 \(\rho_0\in(0,1)\) 使任意跨簇点对都满足
\[
\rho(a,b)\le \rho_0
\qquad
(a\in A_r,\ b\in A_s,\ r\neq s),
\]
则广义缺陷熵
\[
S_{\mathrm{gen}}(A):=-\log\det\mathsf P(A)
\]
满足严格下界
\[
S_{\mathrm{gen}}(A)
\ge
\sum_{r=1}^{R}S_{\mathrm{gen}}(A_r)
+2(-\log\rho_0)\sum_{1\le r<s\le R}\kappa_r\kappa_s.
\]
等价地，
\[
S_{\mathrm{gen}}(A)
\ge
\sum_{r=1}^{R}S_{\mathrm{gen}}(A_r)
+(-\log\rho_0)\left(\kappa^2-\sum_{r=1}^{R}\kappa_r^2\right).
\]
\end{theorem}
\begin{proof}
由结论 \ref{con:xi-terminal-pick-poisson-green-mirror-energy}，
\[
S_{\mathrm{gen}}(A)
=
\sum_{a\in A}(-\log p_a)
+2\sum_{\{a,b\}\subset A}(-\log\rho(a,b)).
\]
将右端按簇内对与跨簇对分拆，得到
\[
S_{\mathrm{gen}}(A)
=
\sum_{r=1}^{R}S_{\mathrm{gen}}(A_r)
+2\sum_{1\le r<s\le R}\ \sum_{a\in A_r,\ b\in A_s}(-\log\rho(a,b)).
\]
跨簇假设给出
\[
-\log\rho(a,b)\ge -\log\rho_0,
\]
因此
\[
S_{\mathrm{gen}}(A)
\ge
\sum_{r=1}^{R}S_{\mathrm{gen}}(A_r)
+2(-\log\rho_0)\sum_{1\le r<s\le R}\kappa_r\kappa_s.
\]
最后用恒等式
\[
2\sum_{r<s}\kappa_r\kappa_s=\kappa^2-\sum_{r=1}^{R}\kappa_r^2
\]
即可改写成第二式。证毕。
\end{proof}

\begin{corollary}[固定总缺陷数与簇数下的强制互能整数极值]\label{cor:derived-pick-poisson-clustered-forced-interaction-extrema}
在定理 \ref{thm:derived-pick-poisson-clustered-quadratic-coercive-lower-bound} 的设定下，进一步固定总缺陷数 \(\kappa\) 与簇数 \(R\)。定义
\[
\mathcal I(\kappa_1,\dots,\kappa_R)
:=
2\sum_{r<s}\kappa_r\kappa_s
=
\kappa^2-\sum_{r=1}^{R}\kappa_r^2.
\]
则：
\begin{enumerate}
  \item \(\mathcal I\) 的最大值由尽可能均衡的整数分簇实现。若
  \[
  \kappa=Rq+s,
  \qquad
  0\le s<R,
  \]
  则
  \[
  \mathcal I_{\max}
  =
  \kappa^2-\bigl((R-s)q^2+s(q+1)^2\bigr);
  \]
  \item \(\mathcal I\) 的最小值由单峰分簇
  \[
  (\kappa-R+1,1,\dots,1)
  \]
  实现，其值为
  \[
  \mathcal I_{\min}=(R-1)(2\kappa-R).
  \]
\end{enumerate}
因此，在固定 \((\kappa,R,\rho_0)\) 下，定理 \ref{thm:derived-pick-poisson-clustered-quadratic-coercive-lower-bound} 的强制互能项具有一个完全显式的整数极值包络。
\end{corollary}
\begin{proof}
由
\[
\mathcal I(\kappa_1,\dots,\kappa_R)=\kappa^2-\sum_{r=1}^{R}\kappa_r^2,
\]
在固定 \(\kappa\) 时，极值问题等价于 \(\sum_r\kappa_r^2\) 的极值问题。

对最大值，凸函数 \(x^2\) 在固定整数和下于分量尽可能均衡时使平方和最小，因此得到所示均衡分簇公式。
对最小值，同一凸性表明平方和在极端偏斜时最大；在约束 \(\kappa_r\ge 1\) 下，这恰由
\[
(\kappa-R+1,1,\dots,1)
\]
实现。直接代入得
\[
\mathcal I_{\min}
=
\kappa^2-\bigl((\kappa-R+1)^2+(R-1)\bigr)
=(R-1)(2\kappa-R).
\]
证毕。
\end{proof}

\begin{theorem}[固定端点通量与统一分离阈值下的二次塌缩律]\label{thm:derived-pick-poisson-fixed-flux-uniform-separation-quadratic-collapse}
设 \(A=\{a_1,\dots,a_\kappa\}\subset\mathbb D\) 为有限点状缺陷，端点 Poisson 权重记为
\[
p_j:=P_{a_j}(-1),
\qquad
p_\Sigma:=\sum_{j=1}^{\kappa}p_j.
\]
若存在 \(\rho_0\in(0,1)\) 使
\[
\rho(a_j,a_k)\le \rho_0
\qquad
(j\neq k),
\]
则有严格不等式
\[
S_{\mathrm{gen}}(A)
\ge
\kappa\log\kappa-\kappa\log p_\Sigma+\kappa(\kappa-1)(-\log\rho_0).
\]
等价地，
\[
\det\mathsf P(A)\le \left(\frac{p_\Sigma}{\kappa}\right)^\kappa \rho_0^{\kappa(\kappa-1)}.
\]
因此，一旦所有缺陷在伪双曲意义下共同靠近，负谱体积不只是线性衰减，而是以 \(\kappa^2\) 级指数塌缩。
\end{theorem}
\begin{proof}
由结论 \ref{con:xi-terminal-pick-poisson-green-mirror-energy}，
\[
S_{\mathrm{gen}}(A)
=
\sum_{j=1}^{\kappa}(-\log p_j)
+2\sum_{1\le j<k\le\kappa}(-\log\rho_{jk}).
\]
其中 \(\rho_{jk}:=\rho(a_j,a_k)\)。

对第一项，由推论 \ref{cor:xi-pick-poisson-energy-fixed-flux-extremal} 的 AM--GM 机制，
\[
\prod_{j=1}^{\kappa}p_j\le \left(\frac{p_\Sigma}{\kappa}\right)^\kappa,
\]
故
\[
\sum_{j=1}^{\kappa}(-\log p_j)\ge \kappa\log\kappa-\kappa\log p_\Sigma.
\]
对第二项，由 \(\rho_{jk}\le \rho_0\) 得
\[
2\sum_{j<k}(-\log\rho_{jk})
\ge
2\binom{\kappa}{2}(-\log\rho_0)
=
\kappa(\kappa-1)(-\log\rho_0).
\]
两式相加即得所示下界；指数化即得行列式上界。证毕。
\end{proof}

\begin{corollary}[有界平均端点通量与统一分离阈值下不存在有限能热力学极限]\label{cor:derived-pick-poisson-no-finite-energy-thermodynamic-limit}
设 \((A_\kappa)_{\kappa\ge 1}\) 为一列缺陷配置，满足：
\begin{enumerate}
  \item 总端点通量至多线性增长，
  \[
  p_\Sigma(A_\kappa)\le C\kappa
  \]
  对某个常数 \(C>0\) 成立；
  \item 存在 \(\rho_0\in(0,1)\) 使
  \[
  \rho(a,b)\le \rho_0
  \]
  对 \(A_\kappa\) 中任意不同点对成立。
\end{enumerate}
则
\[
S_{\mathrm{gen}}(A_\kappa)\ge -\kappa\log C+\kappa(\kappa-1)(-\log\rho_0),
\]
从而
\[
\det\mathsf P(A_\kappa)\le C^\kappa\,\rho_0^{\kappa(\kappa-1)}.
\]
特别地，\(\det\mathsf P(A_\kappa)\) 以严格的 \(\kappa^2\)-级指数速度趋于 \(0\)。
\end{corollary}
\begin{proof}
把 \(p_\Sigma(A_\kappa)\le C\kappa\) 代入定理 \ref{thm:derived-pick-poisson-fixed-flux-uniform-separation-quadratic-collapse}：
\[
\kappa\log\kappa-\kappa\log p_\Sigma(A_\kappa)
\ge
\kappa\log\kappa-\kappa\log(C\kappa)
=
-\kappa\log C.
\]
与该定理中的 \(\kappa(\kappa-1)(-\log\rho_0)\) 项合并，即得所示估计。证毕。
\end{proof}

\begin{theorem}[Schur 补账本中的平均瓶颈原则]\label{thm:derived-pick-poisson-schur-ledger-average-bottleneck}
固定点序 \((a_1,\dots,a_\kappa)\)，令 \(\mathsf P^{(m)}\) 为前 \(m\) 个点诱导的主子矩阵，\(\pi_m\) 为相应 Schur 补通量。定义交叉互能
\[
E_{\mathrm{int}}
:=
S_{\mathrm{gen}}(A)-\sum_{m=1}^{\kappa}(-\log p_m)
=
2\sum_{1\le j<k\le\kappa}(-\log\rho(a_j,a_k)).
\]
则存在某个指标 \(m_\ast\in\{1,\dots,\kappa\}\) 使
\[
\log\frac{p_{m_\ast}}{\pi_{m_\ast}}
\ge
\frac{E_{\mathrm{int}}}{\kappa},
\]
等价地，
\[
\pi_{m_\ast}\le p_{m_\ast}\exp\!\left(-\frac{E_{\mathrm{int}}}{\kappa}\right).
\]
若进一步有统一分离阈值 \(\rho(a_j,a_k)\le \rho_0<1\) 对所有 \(j<k\) 成立，则
\[
\pi_{m_\ast}\le p_{m_\ast}\,\rho_0^{\kappa-1}.
\]
\end{theorem}
\begin{proof}
由结论 \ref{con:xi-terminal-pick-poisson-schur-ledger} 与推论 \ref{cor:xi-pick-poisson-schur-one-step-collapse} 所依赖的账本恒等式，
\[
\det\mathsf P=\prod_{m=1}^{\kappa}\pi_m,
\qquad
-\log \pi_m=-\log p_m+2\sum_{j<m}(-\log\rho(a_j,a_m)).
\]
求和可得
\[
E_{\mathrm{int}}=\sum_{m=1}^{\kappa}\log\frac{p_m}{\pi_m}.
\]
因此由平均值原理，存在某个 \(m_\ast\) 满足
\[
\log\frac{p_{m_\ast}}{\pi_{m_\ast}}
\ge
\frac1\kappa\sum_{m=1}^{\kappa}\log\frac{p_m}{\pi_m}
=
\frac{E_{\mathrm{int}}}{\kappa}.
\]
指数化即得第一条。

若再有 \(\rho(a_j,a_k)\le \rho_0\)，则
\[
E_{\mathrm{int}}
=
2\sum_{j<k}(-\log\rho_{jk})
\ge
2\binom{\kappa}{2}(-\log\rho_0)
=
\kappa(\kappa-1)(-\log\rho_0).
\]
代回第一条得到
\[
\pi_{m_\ast}
\le
p_{m_\ast}\exp\bigl(-(\kappa-1)(-\log\rho_0)\bigr)
=
p_{m_\ast}\rho_0^{\kappa-1}.
\]
证毕。
\end{proof}

\begin{conjecture}[固定总端点通量下的等权--等距极小化]\label{conj:derived-pick-poisson-equalweight-equidistance-minimizer}
固定 \(\kappa\) 与总端点通量 \(p_\Sigma\)。在所有满足
\[
\sum_{j=1}^{\kappa}p_j=p_\Sigma
\]
的有限缺陷配置 \(A=\{a_1,\dots,a_\kappa\}\) 中，广义缺陷熵
\[
S_{\mathrm{gen}}(A)
=
\sum_{j=1}^{\kappa}(-\log p_j)
+2\sum_{j<k}(-\log\rho_{jk})
\]
的极小化应由以下两类对称性同时逼近：
\begin{enumerate}
  \item 端点权重近似等分，
  \[
  p_j\approx p_\Sigma/\kappa;
  \]
  \item 点间伪双曲距离尽可能一致，
  \[
  \rho_{jk}\approx \rho_\ast(\kappa,p_\Sigma)
  \qquad (j\neq k).
  \]
\end{enumerate}
若该猜想成立，则定理 \ref{thm:derived-pick-poisson-fixed-flux-uniform-separation-quadratic-collapse} 中的 AM--GM 自能下界与统一分离阈值互能下界将拼接为一条近乎尖锐的变分律：自能最优要求等权，而互能最优要求等距。
\end{conjecture}
